\documentclass[11pt]{article}
\usepackage[margin=1in]{geometry}
\usepackage{amsmath, amssymb, amsthm}
\usepackage{booktabs}
\usepackage{hyperref}
\usepackage{xcolor}
\usepackage{graphicx}
\usepackage{enumitem}

\title{MatchFormer Fine-Tuning: Self-Supervised Epipolar Loss \\ and Training Setup}
\author{Phase 2 --- Pose-only Supervision}
\date{\today}

\begin{document}
\maketitle

\section{Overview}

We fine-tune MatchFormer (\texttt{litela} variant) using a \emph{fully
self-supervised} objective that requires only camera intrinsics $K$ and
relative pose $(R, t)$ between image pairs. No depth, no ground-truth
correspondences, and no manually annotated matches are used at any point
in the loss computation.

The loss is composed of two complementary epipolar terms on the coarse
confidence matrix: a \textbf{search-reward} term (focal loss with an
epipolar-band pseudo-GT) and a \textbf{precision-steering} term
(differentiable soft Sampson distance). The fine-level loss is disabled
because the only GT fine-offset targets would require depth
reprojection, which we deliberately exclude.

\section{Notation}

\begin{tabular}{ll}
$B$ & batch size \\
$H_c, W_c$ & coarse feature-map height / width (\num{60} / \num{80} at stride \num{8}) \\
$L = H_c W_c$ & number of coarse cells in image~0 ($=\num{4800}$) \\
$S = H_c W_c$ & number of coarse cells in image~1 ($=\num{4800}$) \\
$C \in \mathbb{R}^{B\times L \times S}$ & dual-softmax confidence matrix \\
$p_0(i) \in \mathbb{R}^2$ & pixel-space center of coarse cell $i$ in image~0 \\
$p_1(j) \in \mathbb{R}^2$ & pixel-space center of coarse cell $j$ in image~1 \\
$K \in \mathbb{R}^{3\times 3}$ & shared intrinsics (both images resized to $640\times 480$) \\
$T_0, T_1 \in \mathrm{SE}(3)$ & camera-to-world poses of images 0 and 1 \\
$F \in \mathbb{R}^{3\times 3}$ & fundamental matrix from image~0 to image~1 \\
\end{tabular}

\paragraph{Fundamental matrix from pose + intrinsics.}
Given the relative transform $T_{0\to 1} = T_1^{-1} T_0$ with rotation
$R$ and translation $t$, we build
\begin{equation}
E = [t]_\times R, \qquad F = K^{-\top} E K^{-1},
\label{eq:F}
\end{equation}
computed per-batch-element in \texttt{model/supervision.py}.

\section{Total Loss}

\begin{equation}
\mathcal{L}_{\text{total}}
\;=\;
\lambda_c \cdot \mathcal{L}_{\text{coarse}}
\;+\; \underbrace{\lambda_f \cdot \mathcal{L}_{\text{fine}}}_{\lambda_f = 0}
\end{equation}

with $\lambda_c = 1.0$ and $\lambda_f = 0$. The fine term is removed in
the pose-only setting because a meaningful fine target requires the
depth-reprojected sub-pixel offset of a GT match, which is unavailable.

\begin{equation}
\boxed{\;
\mathcal{L}_{\text{coarse}}
\;=\;
(1 - \lambda_{\text{epi}}) \cdot \mathcal{L}_{\text{focal-epi}}
\;+\;
\lambda_{\text{epi}} \cdot \mathcal{L}_{\text{sampson}}
\;}
\qquad (\lambda_{\text{epi}} = 0.7)
\end{equation}

\subsection{Epipolar mask (pseudo-GT)}

For every pair of coarse grid cells $(i, j)$, define the
point-to-epipolar-line distance in image~1:
\begin{equation}
\ell_i = F \tilde p_0(i),\qquad
d_{ij} \;=\; \frac{\bigl|\tilde p_1(j)^{\top}\,\ell_i\bigr|}{\sqrt{\ell_{i,1}^2 + \ell_{i,2}^2}},
\end{equation}
where $\tilde p = (p, 1)^\top$ is the homogeneous lift. The
\emph{binary} epipolar mask is
\begin{equation}
M_{ij} \;=\; \mathbb{1}\!\left[\, d_{ij} < \tau_{\text{epi}} \,\right],
\qquad \tau_{\text{epi}} = 2 \text{ px}.
\end{equation}
In words: $M_{ij} = 1$ iff image-1 cell $j$ lies within $2$ px of the
epipolar line generated by image-0 cell $i$.

\subsection{$\mathcal{L}_{\text{focal-epi}}$: search reward}

A binary focal loss treating $M$ as soft pseudo-GT:
\begin{equation}
\mathcal{L}_{\text{focal-epi}}
\;=\;
\frac{1}{\sum_{i,j} M_{ij}}
\sum_{i,j}
\Bigl[
-\alpha (1 - C_{ij})^\gamma \log C_{ij}\cdot M_{ij}
\;-\;
(1 - \alpha)\,C_{ij}^{\gamma}\log(1 - C_{ij})\cdot(1 - M_{ij})
\Bigr],
\end{equation}
with $\alpha = 0.25$, $\gamma = 2.0$, and $C$ clipped to
$[10^{-7},\, 1 - 10^{-7}]$ for numerical stability.
Intuitively: \emph{pull confidence onto the epipolar band, push it off
everywhere else}.

\subsection{$\mathcal{L}_{\text{sampson}}$: precision steering}

We compute a \emph{differentiable} Sampson distance via a soft match.

\paragraph{Step A --- Soft match.} For each query $i$, take the
confidence-weighted mean position in image~1:
\begin{equation}
\tilde p_1^{\text{soft}}(i)
\;=\;
\frac{\sum_j C_{ij}\, p_1(j)}{\sum_j C_{ij} + \varepsilon}.
\end{equation}

\paragraph{Step B --- Sampson distance in pixels.} With homogeneous
coordinates $\tilde p_0(i), \tilde p_1^{\text{soft}}(i) \in \mathbb{R}^3$:
\begin{equation}
d_S^2(i)
\;=\;
\frac{\bigl(\tilde p_1^{\text{soft}}(i)^{\top} F\,\tilde p_0(i)\bigr)^2}
     {\|F \tilde p_0(i)\|_{1:2}^2 + \|F^{\top}\tilde p_1^{\text{soft}}(i)\|_{1:2}^2 + \varepsilon}.
\end{equation}

We apply a \textbf{square root} to obtain an $L_1$/Huber-like penalty in
pixel units, which prevents large outliers from dominating the
gradient:
\begin{equation}
d_S(i) \;=\; \sqrt{\,d_S^2(i) + \varepsilon\,}.
\end{equation}

\paragraph{Step C --- Geometric dead zone.} Matches already within
$m = 1$~px of their epipolar line incur zero penalty:
\begin{equation}
\bar d_S(i) \;=\; \max\!\bigl(0,\; d_S(i) - m\bigr).
\end{equation}

\paragraph{Step D --- Confidence-weighted normalization.} With
$C_{\max}(i) = \max_j C_{ij}$ as the row-wise confidence:
\begin{equation}
\mathcal{L}_{\text{sampson}}
\;=\;
\frac{\sum_i C_{\max}(i)\cdot \bar d_S(i)}
     {\sum_i C_{\max}(i) + \varepsilon}.
\end{equation}

The combination of $\sqrt{\cdot}$, the 1-px dead zone, and the
confidence-weighted normalization was chosen to avoid the
\emph{confidence-collapse} failure mode we observed with an earlier
squared-distance formulation.

\subsection{Comparison with SCENES/CAPS}

SCENES and CAPS use a fully per-pair formulation:
\begin{equation}
\mathcal{L}_{\text{SCENES}}
\;=\;
\frac{\sum_{i,j} C_{ij}\cdot d_S(i,j)}{\sum_{i,j} C_{ij}},
\end{equation}
computing $L \cdot S \approx 23$M Sampson distances per image pair. Our
soft-match variant collapses the $j$ dimension first, reducing compute
by $\sim$$S$$\times$ and producing smoother gradients (the soft match
pulls a weighted mean position toward the epipolar line rather than
reweighting discrete candidates). The two losses agree in the
sharp-confidence limit but differ elsewhere, since Sampson is quadratic
in positions.

\section{Why no fine loss}

The fine-level loss
$\mathcal{L}_{\text{fine}} \propto \|\hat o - o_{\text{GT}}\|^2$
requires $o_{\text{GT}}$, the sub-pixel offset of a GT match inside the
coarse window. Recovering $o_{\text{GT}}$ requires either (a) manual
correspondences or (b) depth-based reprojection of the GT coarse point
--- both of which violate the pose-only constraint. We therefore set
$\lambda_f = 0$ and do not call \texttt{fine\_loss} in the training
forward.

Note that the pretrained fine head is retained and still produces
sub-pixel refined matches at inference time; we simply do not update its
weights or supervise its output during fine-tuning.

\section{Self-supervision pipeline}

For every training batch, \texttt{model/supervision.py} populates:
\begin{itemize}[leftmargin=*]
  \item \texttt{data['F\_list']}: list of $B$ fundamental matrices, each
        computed from $K$, $T_0$, $T_1$ via Eq.~\eqref{eq:F}.
  \item \texttt{data['epi\_mask']}: binary tensor of shape $[B, L, S]$,
        entry $(i,j)$ = 1 iff the image-1 cell $j$ is within
        $\tau_{\text{epi}} = 2$~px of the epipolar line of image-0 cell $i$.
  \item \texttt{data['spv\_\{b,i,j\}\_ids']}: empty tensors, retained
        only for API compatibility with \texttt{coarse\_matching.py}.
\end{itemize}

Crucially, \texttt{depth0} is \emph{never} read. The dataset may still
load it for other purposes, but the loss pipeline does not consume it.

\section{Trainable parameters (freeze strategy)}

The loss alone cannot prevent catastrophic forgetting on small
fine-tuning datasets, so we freeze most of the backbone. The current
configuration unfreezes:

\begin{center}
\begin{tabular}{lll}
\toprule
Module & Resolution & Role \\
\midrule
\texttt{backbone.AttentionBlock3} & 1/16 & self + 2$\times$ cross-attention \\
\texttt{backbone.AttentionBlock4} & 1/32 & self + 2$\times$ cross-attention \\
\texttt{backbone.layer1\_outconv}  & 1/4  & FPN projection for fine features \\
\texttt{backbone.layer1\_outconv2} & 1/4  & FPN fusion for fine features \\
\bottomrule
\end{tabular}
\end{center}

All other parameters --- \texttt{AttentionBlock1/2}, the remaining FPN
layers (\texttt{layer2/3/4\_outconv$*$}), and the fine matching head ---
are frozen with \texttt{requires\_grad = False}.

\paragraph{BatchNorm policy.} All BN layers are forced to
\texttt{.eval()} mode in both the initialization path \emph{and} each
\texttt{training\_step}. Their \texttt{weight} and \texttt{bias} have
\texttt{requires\_grad = False}, and \texttt{running\_mean} /
\texttt{running\_var} are frozen to the pretrained statistics. This
prevents BN EMA updates from corrupting the pretrained feature
distribution during small-batch fine-tuning.

\section{Hyperparameters}

\begin{center}
\begin{tabular}{lll}
\toprule
Symbol & Default & Source / CLI flag \\
\midrule
$\lambda_c$           & $1.0$   & \texttt{--lambda\_c} \\
$\lambda_f$           & $0.0$   & \texttt{--lambda\_f} (pose-only: disabled) \\
$\lambda_{\text{epi}}$ & $0.7$   & \texttt{--lambda\_epi} \\
$\alpha$ (focal)      & $0.25$  & hardcoded in \texttt{epi\_focal\_loss} \\
$\gamma$ (focal)      & $2.0$   & hardcoded in \texttt{epi\_focal\_loss} \\
$m$ (Sampson margin)  & $1.0$ px & \texttt{--sampson\_margin} \\
$\tau_{\text{epi}}$   & $2.0$ px & \texttt{--epi\_thresh} \\
dual-softmax temperature & $0.1$ & \texttt{DSMAX\_TEMPERATURE} \\
coarse threshold (train) & $0.2$ & set inside matcher \\
coarse threshold (eval)  & $\{0.1, 0.05, 0.01\}$ & \texttt{--bench\_thresholds} \\
\midrule
learning rate         & $3\times 10^{-5}$ & \texttt{--lr} \\
weight decay          & $1\times 10^{-4}$ & \texttt{WEIGHT\_DECAY} \\
optimizer             & AdamW & \texttt{configure\_optimizers} \\
scheduler             & CosineAnnealingLR & $T_{\max} = $ \texttt{--steps} \\
$\eta_{\min}$         & $1\times 10^{-6}$ & \texttt{--eta\_min} \\
gradient clip         & $1.0$ & \texttt{Trainer.gradient\_clip\_val} \\
batch size            & $8$     & \texttt{--batch} \\
precision             & bf16    & \texttt{--precision} \\
\bottomrule
\end{tabular}
\end{center}

\section{Training data and split}

We use \texttt{ScanNetSimpleDataset} which returns paired RGB images
with their poses and intrinsics. We enforce a per-scene \num{90}/\num{10}
\emph{random} train/test split, seeded by \texttt{--split\_seed 42} so
that local and Colab runs share the same split. Only the \num{90}\%
training portion produces gradients; the held-out \num{10}\% is reserved
for validation and for the in-training pose-AUC benchmark
(Section~\ref{sec:eval}).

Image pairs are sampled with a fixed frame gap of \num{20} frames. Both
images in a pair are resized to $640 \times 480$.

\section{In-training evaluation}
\label{sec:eval}

A \texttt{PoseAUCCallback} fires every \num{1000} training steps and
evaluates the current model on the held-out \num{10}\% split by:
\begin{enumerate}[leftmargin=*]
  \item running a single forward pass at the lowest confidence
        threshold ($0.01$);
  \item post-filtering the predicted matches at each additional
        threshold ($0.05$, $0.10$);
  \item estimating the essential matrix via RANSAC
        (\texttt{thresh}$= 0.5$ px, \texttt{conf}$=0.99999$);
  \item computing rotation and translation errors and aggregating them
        into $\mathrm{AUC}@5^\circ$, $\mathrm{AUC}@10^\circ$, and
        $\mathrm{AUC}@20^\circ$ via the standard cumulative error
        integration.
\end{enumerate}

All metrics are logged to Weights \& Biases under tags
\texttt{bench\_thr\{0.10, 0.05, 0.01\}/\{auc5, auc10, auc20, precision, num\_matches\}}.

\section{Collapse diagnostics}

In addition to the loss decomposition
(\texttt{train/loss\_focal}, \texttt{train/loss\_sampson}), we log the
following health metrics to detect confidence collapse:
\begin{itemize}[leftmargin=*]
  \item \texttt{train/conf\_max}: maximum value of $C$ over the batch.
        Should stay well above $0.1$; a value drifting toward $10^{-3}$
        is a clear collapse signal.
  \item \texttt{train/conf\_mean}: mean of $C$. Serves as a baseline.
  \item \texttt{train/epi\_mask\_frac}: mean of the binary
        \texttt{epi\_mask}. For a well-conditioned pair this is
        typically $\sim$$1$--$3\%$; a value of exactly $0$ indicates
        pose/intrinsic parsing failure or a degenerate pair.
  \item \texttt{train/num\_matches}: count of predicted matches
        surviving mutual-NN + threshold. A drop to single digits
        indicates collapse.
\end{itemize}

\section{Summary}

The fine-tuning objective is a single weighted sum of two
pose-derivable coarse losses:
\[
\mathcal{L}_{\text{total}}
\;=\;
\underbrace{0.3}_{1 - \lambda_{\text{epi}}} \cdot \mathcal{L}_{\text{focal-epi}}
\;+\;
\underbrace{0.7}_{\lambda_{\text{epi}}} \cdot \mathcal{L}_{\text{sampson}}.
\]
No depth, no manual correspondences, and no trainable parameters
outside the two attention blocks and the fine FPN head. Everything else
is a regularizer (BN freeze, cosine LR schedule, gradient clipping) or
an observability aid (collapse diagnostics, in-training pose-AUC
benchmark).

\end{document}
